\bookchapter{Chapitre 35 : Les enfants et le smartphone — Guide éducatif à l'ère du numérique}

\chapterintro{Ce chapitre ne traite pas le smartphone de « monstre » et ne prône pas non plus « plus tôt on commence, plus on est intelligent ». Nous procédons par tranches d'âge : donner des appareils, fixer des règles, enseigner l'éducation numérique, protéger la santé, étape par étape.}

\booksection{I. À quel âge donner un smartphone à son enfant ?}

Il n'y a pas de réponse standard, mais il existe une base d'âge. Trop tôt, le cortex préfrontal de l'enfant n'est pas assez mature, son autocontrôle ne suit pas le « design d'attraction » des algorithmes ; trop tard, il risque d'être « socialement déconnecté » de ses pairs. Voici les recommandations par tranche d'âge :

\begin{itemize}
  \item \textbf{3-6 ans}: Évitez le smartphone autant que possible. En cas d'usage, pas plus de 15 minutes par session, principalement des dessins animés éducatifs sélectionnés par les parents, et obligatoirement\textbf{visionnage accompagné}. À ce stade, l'écran a le plus grand impact sur la vue et le développement de l'attention.
  \item \textbf{7-9 ans}: Peut « emprunter » l'appareil des parents pour les cours en ligne, les histoires audio et les recherches. Commencez à établir la règle « demander avant, rendre après », sans smartphone personnel.
  \item \textbf{10-12 ans}: Envisagez un téléphone « fonctionnel » ou une montre pour enfant, pour les contacts aller-retour école. Si un smartphone est vraiment nécessaire, activez un contrôle parental strict.
  \item \textbf{13 ans et plus}: Au collège, les besoins sociaux augmentent, un smartphone est possible, mais avec un « contrat d'utilisation » et des revues régulières.
\end{itemize}

\begin{tcbinfo}{Principe « assez, c'est assez »}
Le premier smartphone d'un enfant n'a pas besoin d'être un modèle phare. Les performances « excédentaires » du processeur servent souvent à jouer. Un smartphone à environ 1000 yuans avec une batterie suffisante, des performances moyennes et un écran protecteur des yeux, ou un reconditionné officiel, facilite au contraire le contrôle. Budget de référence : 1500-2500 yuans suffisent pour la plupart des élèves.

\end{tcbinfo}

\booksubsection{La montre intelligente comme solution transitoire}

Si vous voulez juste résoudre le problème « pouvoir joindre l'enfant après l'école », une montre pour enfant est plus adaptée qu'un smartphone. Elle ne peut que téléphoner, envoyer des messages vocaux et localiser, sans navigateur ni courtes vidéos, limitant naturellement les « portes d'entrée vers l'addiction » :

\begin{itemize}
  \item \textbf{Montre XTC (Xiao Tian Cai)}: L'écosystème le plus complet, la fonction sociale « touche pour ajouter un ami » entre pairs est son argument principal, mais l'écosystème d'applications dédié est relativement fermé.
  \item \textbf{Montre pour enfant Huawei}: Bonne interaction avec les smartphones Huawei, localisation précise, prise en charge des appels vidéo, bon rapport qualité-prix.
  \item \textbf{Montre pour enfant Xiaomi}: Prix abordable, fonctions de base complètes, adaptée aux familles au budget limité.
\end{itemize}

Lors de l'achat d'une montre, choisissez un modèle prenant en charge\textbf{la 4G toutes bandes}(la 2G/3G est progressivement retirée), et vérifiez les tarifs des forfaits enfants.

\booksection{II. Guide complet du contrôle parental}

Avant de donner l'appareil à l'enfant, la première chose n'est pas de l'allumer, mais de configurer le contrôle parental. Les deux plateformes ont des solutions natives :

\booksubsection{iOS : Temps d'écran}

\begin{itemize}
  \item \textbf{Limites d'utilisation}: Réglages > Temps d'écran > Limites d'utilisation, permet de définir une durée quotidienne par catégorie d'application (par exemple 30 minutes pour les jeux).
  \item \textbf{Temps d'arrêt}: Définissez un « couvre-feu » quotidien (par exemple 22h00-7h00), pendant lequel toutes les applications sont désactivées, seul le téléphone reste disponible.
  \item \textbf{Restrictions de contenu}: Filtrez les films, séries, pages web et installations d'applications par tranche d'âge, et interdisez les achats intégrés.
  \item \textbf{Partage familial}: Gérez à distance l'appareil de l'enfant depuis l'iPhone des parents, approuvez les demandes de téléchargement, consultez les rapports d'utilisation.
\end{itemize}

\booksubsection{Android : Google Family Link / contrôle parental de marque}

\begin{itemize}
  \item \textbf{Google Family Link}: Solution officielle de Google, gestion à distance de la durée d'utilisation, verrouillage de l'appareil, consultation des rapports d'activité. La disponibilité sur les Android nationaux dépend du modèle.
  \item \textbf{Espace enfant Xiaomi}: Réglages > Espace enfant, permet de définir la durée d'utilisation, les applications autorisées, la sortie nécessite le mot de passe parental.
  \item \textbf{Mode étudiant Huawei}: Réglages > Utilisation saine > Mode étudiant, limite les plages horaires et les applications, prend en charge la gestion à distance.
  \item \textbf{Mode enfant OPPO/vivo}: S'active dans « Réglages » ou l'application « Espace enfant » dédiée, principe similaire.
\end{itemize}

\begin{center}
\begin{tabularx}{\linewidth}{|X|X|X|}
\hline
\textbf{Fonction} & \textbf{iOS Temps d'écran} & \textbf{Android Family Link / solution de marque} \tabularnewline
\hline
Limites d'utilisation & Pris en charge, par catégorie & Pris en charge, par application \tabularnewline
Temps d'arrêt & Pris en charge & Pris en charge (varie selon la marque) \tabularnewline
Filtrage de contenu & Au niveau système, assez complet & Dépend de la marque, inégal \tabularnewline
Gestion à distance & Partage familial & Family Link / cloud de marque \tabularnewline
Approbation d'installation & Approbation parentale requise & Family Link le prend en charge \tabularnewline
\hline
\end{tabularx}
\end{center}

\begin{tcbwarn}{Le contrôle parental n'est pas « une fois pour toutes »}
Les enfants trouvent toujours des moyens de contourner (les tutoriels « comment contourner le temps d'écran » abondent en ligne). La technique n'est qu'une ligne de fond, la vraie défense est la communication et la confiance parent-enfant. Il est recommandé de revoir le rapport d'utilisation avec l'enfant chaque mois, plutôt que de surveiller unilatéralement.

\end{tcbwarn}

\booksection{III. Mode adolescent}

Les principales plateformes de contenu intègrent un « mode adolescent » qui, une fois activé, limite la durée d'utilisation, filtre les contenus inappropriés et désactive les paiements et pourboires. C'est une fonction obligatoire imposée par la réglementation nationale, mais\textbf{les parents doivent l'activer et définir un mot de passe}:

\begin{itemize}
  \item \textbf{Mode adolescent Douyin/Kuaishou}: 40 minutes par jour (interdit de 22h00 à 6h00), seuls les contenus scientifiques et éducatifs sélectionnés sont affichés, pas de live, pas de pourboires, pas de recharge.
  \item \textbf{Mode adolescent WeChat}: Limite WeChat Pay (nécessite une autorisation parentale), restreint le contenu des comptes publics et des chaînes vidéo, les fonctions « Secouer » et « À proximité » sont désactivées par défaut.
  \item \textbf{Système anti-addiction aux jeux}: Authentification par nom réel et limitation de temps au niveau national. Les mineurs ne peuvent se connecter aux jeux que le vendredi, samedi, dimanche et jours fériés de 20h00 à 21h00, 1 heure par jour.
\end{itemize}

\textbf{N'utilisez pas la pièce d'identité des parents pour « emprunter un compte »}\par

Beaucoup d'enfants utilisent la pièce d'identité de leurs parents pour enregistrer un compte de jeu et contourner l'anti-addiction, ce qui enfreint le contrat utilisateur et prive l'enfant de protection. Les parents doivent garder leurs mots de passe de paiement et informations d'identité, et ne pas « prêter de compte » par commodité.

\booksection{IV. Éducation à la citoyenneté numérique}

La technologie peut bloquer 90 \% des risques, les 10 \% restants dépendent du jugement de l'enfant. Plutôt que de dire « interdiction d'aller sur Internet », il est plus important de leur apprendre « comment y aller » :

\begin{itemize}
  \item \textbf{Protection de la vie privée}: Ne pas révéler son vrai nom, son école ou son adresse à des inconnus ; ne pas envoyer de photos personnelles, de pièce d'identité ou de photos de famille ; la localisation est désactivée par défaut, activée uniquement en cas de besoin.
  \item \textbf{Reconnaissance et gestion du cyberharcèlement}: Insultes, isolement, retouches malveillantes de photos sont du harcèlement. Dites à l'enfant de ne pas répondre, de ne pas supprimer les preuves, et d'en parler immédiatement à un parent. Sur WeChat, un appui long sur le message permet de signaler, et la plateforme sanctionnera le compte harceleur.
  \item \textbf{Esprit critique sur le contenu}: Apprenez à l'enfant à distinguer « fait » et « opinion », à repérer les titres putaclic et les styles sensationnalistes ; à comprendre que l'algorithme ne propose que ce que vous voulez voir, ce qui ne représente pas le monde entier ; à ne jamais croire les publicités « repas gratuit ».
  \item \textbf{Sécurité en ligne}: Ne pas cliquer sur des liens inconnus, ne pas scanner de QR codes d'origine inconnue, ne pas télécharger d'applications hors des boutiques officielles ; ne partager son mot de passe avec personne (pas même un « meilleur ami ») ; toute demande de code de vérification est considérée comme une arnaque.
\end{itemize}

\begin{tcbpractice}{Exercice : signez un « contrat d'utilisation du smartphone » avec votre enfant}
Écrire les règles, les signer ensemble et les afficher bien en vue est bien plus efficace qu'un accord verbal. Suggestions :

\begin{itemize}
  \item Plages d'utilisation : par exemple « disponible de 19h30 à 20h30 après les devoirs ».
  \item Zones interdites : à table, au lit avant de dormir, en classe, le smartphone est rangé.
  \item Limites de contenu : ne pas installer d'applications sans en informer les parents, ne pas discuter en privé avec des inconnus.
  \item Conséquences en cas d'infraction : par exemple « un dépassement = durée réduite de moitié le lendemain ».
  \item Engagement parental : ne pas « scroller » devant l'enfant au-delà de la durée convenue, donner l'exemple.
\end{itemize}

\end{tcbpractice}

\booksection{V. Habitudes d'utilisation saines}

Les yeux, la colonne cervicale et le sommeil des enfants sont en pleine croissance, l'impact du smartphone sur leur corps est plus important que sur les adultes :

\begin{itemize}
  \item \textbf{Recommandations de temps d'écran}: 3-6 ans : pas plus d'1 heure par jour ; 6-12 ans : pas de limite pour l'apprentissage, 1 heure maximum pour le divertissement ; 12 ans et plus : à négocier, mais 2 heures maximum pour le divertissement.
  \item \textbf{Habitudes de protection des yeux}: Respectez la « règle 20-20-20 » — toutes les 20 minutes, regardez à 20 pieds (environ 6 mètres) pendant 20 secondes ; pas de smartphone 1 heure avant le coucher, la lumière bleue supprime la mélatonine et affecte le sommeil.
  \item \textbf{Gestion de la posture}: L'écran du smartphone à hauteur des yeux, ne pas baisser la tête ni voûter les épaules ; se lever et bouger après 30 minutes d'utilisation continue.
  \item \textbf{Règles familiales communes}: Définissez des « zones sans smartphone » à table, dans la chambre et le bureau ; lors des films en famille et des activités en plein air, tout le monde range son smartphone, pour que « poser le smartphone » devienne une culture familiale et non une punition pour l'enfant.
\end{itemize}

\begin{tcbinfo}{Les parents sont le meilleur « pare-feu »}
Les recherches montrent une forte corrélation entre le temps d'usage du smartphone par les parents et celui des enfants. Exiger que l'enfant ne regarde pas le smartphone tout en scrollant des courtes vidéos jusqu'à tard le soir rend les règles inefficaces. Pour changer l'enfant, commencez par poser votre propre smartphone.

\end{tcbinfo}

\booksection{VI. Patrimoine numérique et enfants}

Cette section est souvent négligée mais importante. Les actifs numériques des enfants (comptes, photos, soldes rechargés, équipements de jeu) posent eux aussi des questions de succession et de confidentialité :

\begin{itemize}
  \item \textbf{Succession des actifs numériques de l'enfant}: En cas de décès malheureux de l'enfant, la succession de ses comptes sociaux, de jeux et de photos dans le cloud suit la même procédure que pour les adultes, nécessitant un acte de décès et un justificatif de lien de parenté. Il est recommandé aux parents de configurer à l'avance un « contact de succession » ou un « contact d'urgence » pour les comptes importants de l'enfant (comme le contact de succession Apple, la succession de compte WeChat).
  \item \textbf{Confidentialité numérique des mineurs}: La « Loi sur la protection des informations personnelles » stipule que le traitement des données personnelles des mineurs de moins de 14 ans nécessite le consentement du tuteur et des règles de confidentialité spécifiques. Les parents doivent examiner régulièrement la politique de confidentialité des applications utilisées par l'enfant, et désactiver les « recommandations personnalisées » et le « partage de données ».
  \item \textbf{Recommandations pour l'enregistrement de comptes}: Les comptes enregistrés pour l'enfant doivent utiliser un e-mail/numéro contrôlable par les parents, pour une gestion unifiée pendant la minorité ; le transfert du compte peut se faire à la majorité.
\end{itemize}

\begin{tcbwarn}{L'autre face du « patrimoine numérique des enfants »}
Les enfants sont « enregistrés » dès la naissance — photos de croissance publiées par les parents, vidéos envoyées par la maternelle, données de santé de la montre connectée. Ces données, une fois téléchargées, sont difficiles à supprimer complètement. Avant de publier, demandez-vous : dans dix ans, cette photo gênera-t-elle l'enfant ? Respecter la personnalité numérique de l'enfant commence par s'abstenir de tout publier.

\end{tcbwarn}

\begin{tcbsummary}{Résumé du chapitre}
\begin{itemize}
  \item L'âge pour donner un smartphone n'a pas de réponse standard, mais 3-6 ans évitent, avant 13 ans transition par montre ou téléphone fonctionnel
  \item « Assez, c'est assez » : le premier smartphone à environ 1000 yuans, pas un phare ; performances excédentaires = tentations excédentaires
  \item Le temps d'écran iOS et Family Link/contrôle parental Android sont deux lignes de défense techniques, plus stables ensemble
  \item Le mode adolescent et l'anti-addiction aux jeux sont obligatoires, mais les parents doivent les activer et protéger le mot de passe
  \item La citoyenneté numérique compte plus que le filtrage technique : apprenez à l'enfant à reconnaître le harcèlement, protéger sa vie privée, critiquer le contenu
  \item Les habitudes saines commencent par les parents — si vous posez le smartphone, l'enfant le fera aussi
  \item Le patrimoine et la confidentialité numériques des enfants sont des sujets émergents ; configurez un « contact de succession » et abstenez-vous de trop publier
\end{itemize}

\end{tcbsummary}
